\newpage

% Переключаем на A3 landscape
\KOMAoptions{paper=A3,paper=landscape}
\areaset{390mm}{270mm}
\fancyheadoffset{0pt} % обновляем хедер и футер

\topmargin=-2.6cm % Команда \topmargin задаёт верхнее поле страницы. При этом поле отсчитывается не от левого края листа, а от линии, параллельной краю листа и отстоящей от него на 1 дюйм
\oddsidemargin=0cm % отступ слева https://zns.susu.ru/IT/latex/Marking/marking.html
\headsep=0cm
\footskip=-0.8cm
\setlength{\headheight}{1cm}


\color{unidarkgreen}{\section[(рекомендуемое) Структурно-функциональная схема устройства]{(рекомендуемое)\\Структурно-функциональная схема устройства}\label{app:fsu}}
\color{black}

% Вставляем рисунки
\newcounter{filecounter}
\setcounter{filecounter}{1}
\loop
    \edef\filename{\apppath/Приложение. ФСУ/_latex/img/\arabic{filecounter}.pdf}
    \IfFileExists{\filename}{%
        \begin{figure}[H]
        \centering
        % Для первой страницы используем высоту 0.9\textheight (запас под заголовок),
        % для остальных – полную \textheight
        \ifnum\value{filecounter}=1
            \resizebox{!}{0.9\textheight}{\includegraphics{\filename}}
        \else
            \resizebox{!}{\textheight}{\includegraphics{\filename}}
        \fi
        \end{figure}
        \clearpage
        \stepcounter{filecounter}
    }{%
        \setcounter{filecounter}{0}
    }
    \ifnum\value{filecounter}>0
\repeat

\clearpage

\KOMAoption{paper}{a4, portrait}
\areaset{267mm}{180mm}
\fancyheadoffset{0pt} % обновляем хедер и футер

\newgeometry{
  top=14mm,
  bottom=11mm,
  right=10mm,
  left=22mm,
}
